\documentclass[11pt,twoside]{amsart}
\usepackage{amssymb, amsmath, enumerate, palatino, hyperref}
\usepackage{caption}
\usepackage{fullpage}
\usepackage{tikz}
\usetikzlibrary{arrows.meta,positioning}

\usepackage[T1]{fontenc}
\renewcommand{\labelitemi}{\guillemotright}
\usepackage{mathrsfs}
\hypersetup{hidelinks}

\theoremstyle{plain}
\newtheorem{prop}{Proposition}
\newtheorem{lemma}[prop]{Lemma}
\newtheorem{thm}[prop]{Theorem}
\newtheorem{cor}[prop]{Corollary}
\theoremstyle{remark}
\newtheorem{rmk}[prop]{Remark}
\newtheorem*{rmk*}{Remark}
\newtheorem{prob}{Problem}
\newtheorem*{bonus*}{Bonus Problem}
\theoremstyle{definition}
\newtheorem{ex}[prop]{Example}
\newtheorem{defn}[prop]{Definition}

\newcommand{\RR}{\mathbb{R}}
\newcommand{\ZZ}{\mathbb{Z}}
\newcommand{\Ztn}{(\mathbb{Z}/2)^{n}}
\newcommand{\tr}{\operatorname{tr}}
\newcommand{\spec}{\operatorname{spec}}
\newcommand{\wt}{\omega}
\newcommand{\kap}{\kappa}
\newcommand{\adj}{\operatorname{adj}}
\newcommand{\Uh}{\widehat U}

\title{Math 372: Combinatorics\\ Homework 03}
\author{Due Friday, 9 October, 10\textsc{p.m.}, via Gradescope}

\begin{document}
\maketitle

\begin{rmk*}
\textbf{When you can start.} Problems 1 and 2 use only what we have done through Monday,
September~28. Problem~3 needs Wednesday the 30th and Friday, October~2; Problem~4(c) and (d)
need Monday, October~5. As before, the set is released early so that the first half can be
done first; do not wait.
\end{rmk*}

\begin{rmk*}
\textbf{Conventions.} Graphs are finite, undirected, and loopless. For a subgroup $G\le S_n$
acting on $B_n$, the quotient $B_n/G$ is the poset of orbits with $o\le o'$ iff some
$x\in o$ lies below some $y\in o'$; for an orbit $o$, $v_o=\sum_{x\in o}x\in\RR B_n$ is its
orbit sum, and $\Uh_i\colon\RR(B_n/G)_i\to\RR(B_n/G)_{i+1}$ is the operator
$\Uh_i(o)=\sum_{o'}c_{oo'}\,o'$ with $c_{oo'}=\#\{x\in o: x<y\}$ for any $y\in o'$, as in
class. For a graph $G$ with vertices $v_1,\dots,v_p$, $L=L(G)$ is the Laplacian (degrees on the
diagonal, $-1$ for each edge, $0$ otherwise), $L_0$ is $L$ with one row and the same column
deleted, and $\kap(G)$ is the number of spanning trees.
\end{rmk*}

\begin{rmk*}
\textbf{On these solutions.} Each problem is worth up to five points for mathematical content
and up to two for writing. Write for another student in this course: say what you are going to
show before you show it, justify each step, and use complete sentences. End your submission with
a short \emph{Methods and tools} note --- three or four sentences on what software you used and
for what, and how you checked anything a machine produced for you.
\end{rmk*}

%%% Problem 1 %%%
\begin{prob}[necklaces]
Let $\sigma=(1\,2\,\cdots\,n)\in S_n$ and $G=\langle\sigma\rangle\cong\ZZ/n$, acting on $B_n$.
An orbit of $G$ on $k$-subsets is a \emph{necklace} with $n$ beads, $k$ of them black.
\begin{enumerate}[(a)]
\item For $n=6$, list the orbits rank by rank, drawing each necklace and recording the size of
each orbit. (You should find rank sizes $1,1,3,4,3,1,1$.)
\item Still for $n=6$, write the matrix of $\Uh_2\colon\RR(B_6/G)_2\to\RR(B_6/G)_3$ in the
orbit bases, computing each coefficient $c_{oo'}$ from its definition, and verify that
$\Uh_2$ is one-to-one.
\item Now let $n=p$ be prime. Show that for $0<k<p$ every orbit of $G$ on $k$-subsets has
exactly $p$ elements. Deduce a formula for the number of necklaces with $p$ beads, $k$ of them
black, and conclude that $p$ divides $\binom pk$ for $0<k<p$.
\item What does the theorem from class say about the largest antichain in $B_6/G$? The
theorem proves that a chain decomposition of $B_6/G$ into symmetric chains exists; find one by
hand, using your Hasse diagram from (a).
\end{enumerate}
\end{prob}

%%% Problem 2 %%%
\begin{prob}[the grid, again]
Fix $1\le k<n$ and let $G=S_k\times S_{n-k}\le S_n$ be the subgroup of permutations that map
$[k]=\{1,\dots,k\}$ to itself. For $0\le a\le k$ and $0\le b\le n-k$ let
\[
O_{a,b}=\bigl\{\,x\in B_n:\ |x\cap[k]|=a\ \text{and}\ |x\smallsetminus[k]|=b\,\bigr\}.
\]
\begin{enumerate}[(a)]
\item Show that the $O_{a,b}$ are exactly the orbits of $G$ on $B_n$, and that
$O_{a,b}\le O_{a',b'}$ in $B_n/G$ if and only if $a\le a'$ and $b\le b'$. Conclude that
$B_n/G$ is isomorphic to the grid poset of Problem~3 on Problem Set~2 (with the roles of
$a,b$ there played by $k,n-k$ here).
\item Compute the coefficients $c_{oo'}$ for $o=O_{a,b}$ and $o'$ an orbit of the next
rank, and write $\Uh_i(O_{a,b})$ as a combination of orbits.
\item Take $k=1$. Write $\Uh_i$ as a $2\times2$ matrix for $1\le i<n/2$ and check directly
that it is one-to-one. What happens for $i\ge n/2$, and why is that consistent with the
theorem?
\item In one paragraph: the theorem from class reproves part (c) of Problem~3 on Problem
Set~2. What did the problem-set proof provide that the theorem does not?
\end{enumerate}
\end{prob}

%%% Problem 3 %%%
\begin{prob}[Laplacians]
\begin{enumerate}[(a)]
\item Write down $L$ for the complete bipartite graph $K_{2,3}$, with the two vertices of the
small side listed first. Compute $\det L_0$ twice, deleting a vertex from each side, and confirm
the answer by counting the spanning trees of $K_{2,3}$ directly.
\item Let $G$ be connected. Prove, \emph{without} using the Matrix-Tree Theorem, that
$\det L_0$ does not depend on which row and column are deleted. (Suggestion: let $\adj L$ be
the adjugate, so that $L\cdot\adj L=\adj L\cdot L=(\det L)\,I=0$. Explain why the kernel of $L$
is spanned by the all-ones vector $\mathbf 1$, deduce that every column and every row of
$\adj L$ is a multiple of $\mathbf 1$, and conclude that $\adj L=c\,J$ for a constant $c$,
where $J$ is the all-ones matrix. Then identify the diagonal entries of $\adj L$.)
\item Let $G$ be disconnected. Show that $\det L_0=0$ for every choice of deleted vertex, so
the Matrix-Tree Theorem remains true, in the form ``$0=0$,'' for disconnected graphs.
\item (\emph{Deletion--contraction.}) For an edge $e$ of $G$, let $G-e$ be $G$ with $e$
removed and $G/e$ be $G$ with $e$ contracted to a single vertex (keeping any multiple edges
that result). Prove that
\[
\kap(G)=\kap(G-e)+\kap(G/e),
\]
and use this to compute $\kap(F_n)$ for $n\le 4$, where the \emph{fan} $F_n$ is a path on
$n$ vertices together with one more vertex joined to every vertex of the path. Guess the
general pattern.
\end{enumerate}
\end{prob}

%%% Problem 4 %%%
\begin{prob}[spanning trees from eigenvalues]
\begin{enumerate}[(a)]
\item Verify Cauchy--Binet by hand for
$A=\bigl(\begin{smallmatrix}1&2&3\\4&5&6\end{smallmatrix}\bigr)$ and $B=A^{\mathsf T}$:
compute $\det(AA^{\mathsf T})$ directly, and as the sum over the three $2$-subsets of columns.
\item Show that the eigenvalues of $L(K_{m,n})$ are $0$ and $m+n$, each once, $m$ with
multiplicity $n-1$, and $n$ with multiplicity $m-1$. (Look for eigenvectors supported on one
side of the bipartition whose entries sum to zero, and for eigenvectors constant on each side.)
Deduce that $\kap(K_{m,n})=m^{\,n-1}n^{\,m-1}$, and check it against Problem~3(a).
\item The cube $Q_n$ is $n$-regular and its adjacency eigenvalues are $n-2i$ with multiplicity
$\binom ni$, $0\le i\le n$. Use the eigenvalue form of the Matrix-Tree Theorem to find
$\kap(Q_n)$ in closed form, and evaluate it for $n=2,3,4$.
\item Let $G_4$ be the two-flip graph of Problem~1 on Problem Set~2, whose spectrum you
computed there. It has two isomorphic connected components. Use the spectrum to find the
number of spanning trees of one component, and say what graph that component is.
\end{enumerate}
\end{prob}

%%% Bonus %%%
\begin{bonus*}[trees through an edge]
Show that the number of spanning trees of $K_n$ that contain a given edge is $2n^{n-3}$.
(Count pairs (tree, edge of the tree) in two ways.) Deduce $\kap(K_n-e)$, and check your
answer for $n=4$ with the Matrix-Tree Theorem.
\end{bonus*}

\end{document}
